\begin{lemma}[The almost-everywhere limit of the complex Riesz images, and its uniqueness in measure]
Let $f \in L^1(\gamma_n;\mathbb{C})$ and let $(v_k)_{k\in\mathbb{N}}$ be a
sequence in $L^2(\gamma_n;\mathbb{C})$.  Write $\iota_{\mathbb{C}}$ for the
inclusion of \noderef{ComplexToL1} and $R_\gamma^{\mathbb{C}}$ for the transform
of \noderef{ComplexRieszTransform}.  Assume the three hypotheses
\begin{enumerate}
  \item[(A)] for all $u,w \in L^2(\gamma_n;\mathbb{C})$ and every
    $\varepsilon > 0$,
    \[
      \gamma_n\bigl(\{x : \varepsilon \leq |R_\gamma^{\mathbb{C}} u(x)
        - R_\gamma^{\mathbb{C}} w(x)|_{\mathbb{C}^n}\}\bigr)
        \leq \frac{16}{\pi}\cdot
          \frac{\|\iota_{\mathbb{C}} u - \iota_{\mathbb{C}} w\|_{L^1}}{\varepsilon} ;
    \]
  \item[(L)] $\iota_{\mathbb{C}} v_k \to f$ in $L^1(\gamma_n;\mathbb{C})$;
  \item[(C)] $\|\iota_{\mathbb{C}} v_j - \iota_{\mathbb{C}} v_{j+1}\|_{L^1}
    \leq 8^{-j}$ for every $j$.
\end{enumerate}
Then there is an element $g$ of $\mathcal{M}(\gamma_n;\mathbb{C}^n)$, the space
of $\gamma_n$-almost-everywhere equivalence classes of measurable
$\mathbb{C}^n$-valued functions on $\mathbb{R}^n$, such that
\begin{enumerate}
  \item $R_\gamma^{\mathbb{C}} v_k \to g$ pointwise $\gamma_n$-a.e.; and
  \item for \emph{every} sequence $(w_k)$ in $L^2(\gamma_n;\mathbb{C})$ with
    $\iota_{\mathbb{C}} w_k \to f$ in $L^1(\gamma_n;\mathbb{C})$, the sequence
    $R_\gamma^{\mathbb{C}} w_k$ converges to $g$ in measure.
\end{enumerate}
\end{lemma}
\begin{proof}
\emph{Part 1: the almost-everywhere limit.}
Put $E_j = \{x : 2^{-j} \leq |R_\gamma^{\mathbb{C}} v_j(x)
 - R_\gamma^{\mathbb{C}} v_{j+1}(x)|_{\mathbb{C}^n}\}$.  By (A) with
$\varepsilon = 2^{-j}$ and then (C),
\[
  \gamma_n(E_j) \leq \frac{16}{\pi}\cdot
    \frac{\|\iota_{\mathbb{C}} v_j - \iota_{\mathbb{C}} v_{j+1}\|_{L^1}}{2^{-j}}
    \leq \frac{16}{\pi}\cdot 2^{j}\cdot 8^{-j}
    = \frac{16}{\pi}\cdot 4^{-j} .
\]
The geometric series $\sum_j \tfrac{16}{\pi}4^{-j}$ converges, so
$\sum_j \gamma_n(E_j) < \infty$ and the Borel--Cantelli lemma (shared context)
gives: for $\gamma_n$-a.e.\ $x$ the set $\{j : x \in E_j\}$ is finite.  Fix such
an $x$ and let $B$ bound that finite set.  For $j \geq B+1$ we have
$|R_\gamma^{\mathbb{C}} v_j(x) - R_\gamma^{\mathbb{C}} v_{j+1}(x)| < 2^{-j}$, so
all but finitely many terms of
$\sum_j |R_\gamma^{\mathbb{C}} v_j(x) - R_\gamma^{\mathbb{C}} v_{j+1}(x)|$ are
dominated by the convergent series $\sum_j 2^{-j}$; hence that series converges,
and a sequence in the complete space $\mathbb{C}^n$ whose consecutive distances
form a convergent series is Cauchy, hence convergent.  Define $g(x)$ to be this
limit at every $x$ where $(R_\gamma^{\mathbb{C}} v_j(x))_j$ converges, and
$g(x) = 0$ elsewhere.  Then $R_\gamma^{\mathbb{C}} v_j \to g$ pointwise
$\gamma_n$-a.e., which is clause (1); and $g$ is measurable, being a
$\gamma_n$-a.e.\ pointwise limit of the measurable functions
$R_\gamma^{\mathbb{C}} v_j$, so it determines an element of
$\mathcal{M}(\gamma_n;\mathbb{C}^n)$.

\emph{Part 2: every approximating sequence has the same limit in measure.}
Since $\gamma_n$ is a finite measure (\noderef{GaussIsProb}), the almost
everywhere convergence of Part 1 implies convergence in measure: for every
$\delta > 0$,
\begin{equation*}
  \gamma_n\bigl(\{x : \delta \leq |R_\gamma^{\mathbb{C}} v_k(x)
    - g(x)|_{\mathbb{C}^n}\}\bigr) \longrightarrow 0 . \tag{E}
\end{equation*}
Let now $(w_k)$ be any sequence in $L^2(\gamma_n;\mathbb{C})$ with
$\iota_{\mathbb{C}} w_k \to f$ in $L^1$, and fix $\varepsilon > 0$.  By the
triangle inequality in $\mathbb{C}^n$, if
$\varepsilon \leq |R_\gamma^{\mathbb{C}} w_k(x) - g(x)|$ then at least one of
$|R_\gamma^{\mathbb{C}} w_k(x) - R_\gamma^{\mathbb{C}} v_k(x)|$ and
$|R_\gamma^{\mathbb{C}} v_k(x) - g(x)|$ is at least $\varepsilon/2$; so
\[
  \{\varepsilon \leq |R_\gamma^{\mathbb{C}} w_k - g|\}
    \subseteq \{\tfrac{\varepsilon}{2}
        \leq |R_\gamma^{\mathbb{C}} w_k - R_\gamma^{\mathbb{C}} v_k|\}
      \cup \{\tfrac{\varepsilon}{2} \leq |R_\gamma^{\mathbb{C}} v_k - g|\} ,
\]
and monotonicity and subadditivity of $\gamma_n$ (shared context) give
\[
  \gamma_n\bigl(\{\varepsilon \leq |R_\gamma^{\mathbb{C}} w_k - g|\}\bigr)
    \leq \gamma_n\bigl(\{\tfrac{\varepsilon}{2}
          \leq |R_\gamma^{\mathbb{C}} w_k - R_\gamma^{\mathbb{C}} v_k|\}\bigr)
      + \gamma_n\bigl(\{\tfrac{\varepsilon}{2}
          \leq |R_\gamma^{\mathbb{C}} v_k - g|\}\bigr) .
\]
The second summand tends to $0$ by (E) with $\delta = \varepsilon/2$.  For the
first, $\|\iota_{\mathbb{C}} w_k - \iota_{\mathbb{C}} v_k\|_{L^1} \to 0$, because
both $\iota_{\mathbb{C}} w_k$ and $\iota_{\mathbb{C}} v_k$ converge to $f$ in
$L^1$ by hypothesis and by (L) respectively, so (A) at the level
$\varepsilon/2$ bounds the first summand by
$32\,\|\iota_{\mathbb{C}} w_k - \iota_{\mathbb{C}} v_k\|_{L^1}/(\pi\varepsilon)
 \to 0$.  Both summands are nonnegative and tend to $0$, so their sum does, and
the left-hand side, squeezed between $0$ and that sum, tends to $0$.  That is
convergence in measure of $R_\gamma^{\mathbb{C}} w_k$ to $g$, which is
clause (2).
\end{proof}
